\begin{textsample}{POS Dim 3 – pa – Score 8.00 – pa/pa\_inf\_9.txt}  \label{ex:f3_pos_279}
3.1.3 O \textbf{Brincar} como Direito Discorrer sobre o “ \textbf{brincar} como direito ” e intercruzar suas bases com a Educação \textbf{Infantil} mostra -se importante primeiro situar o leitor sobre qual lugar está sendo falado, quando se vai ao encontro da ação do \textbf{brincar} e que concepções estão atreladas e caminham juntas a esta ação.
Oliveira ( 2012 ) conceitua o \textbf{brincar} como algo aprendido nas interações sociais e no contato com as manifestações culturais produzidas e destaca a atuação do professor como colaborador na ampliação e redimensionamento da ação do \textbf{brincar} à medida que observa, reflete, planeja e intervém oferecendo às \textbf{crianças} novos elementos disponíveis na cultura para dialogar com as \textbf{crianças} em diferentes espaços e tempos.
Nessa linha de pensamento, assume -se o \textbf{brincar}, no âmbito deste documento, como ato revelador da existência da \textbf{criança} como pessoa, da sua identidade, estando, portanto, atrelado a sua própria razão de ser no mundo, afinal \textbf{brincar} é algo muito sério porque envolve uma gama de conhecimentos complexos e refinados, elaborados e reelaborados, por pessoas humanas ocupantes de um tempo histórico chamado infância.
É basilar que antes de se aprofundar em qualquer análise, se assuma o \textbf{brincar} como direito da \textbf{criança}, justificando tal atitude como uma situação de justiça social, de políticas públicas, de direito universal que : “ em todas as medidas relativas à infância será dada prioridade aos melhores interesses da infância ” ( ONU, 2002, p. 13 ).
Outro movimento fundamental ao qual se é impelido é desarticular a noção do ato de \textbf{brincar} apenas a Educação \textbf{Infantil} ou a escola, tampouco o \textbf{brincar} como campo restrito da ação da \textbf{criança}.
É necessário alargar o olhar para as infâncias e as \textbf{crianças} para melhor compreendermos que o \textbf{brincar} não está amarrado à escolarização, e é imperioso desescolarizá -lo.
O \textbf{brincar} pertence à vida do homem na Terra, mais particularmente ao tempo do ser-criança e se articula com os modos de ser e as produções culturais terrenas.
Ao assumir a concepção de \textbf{criança} como ser de capacidades e potencialidades, compreende -se que elas conseguem, mesmo antes de falar ou andar, elaborar e reelaborar conhecimentos complexos concretamente observáveis nas ações que realizam e nas linguagens que articulam, definindo assim o compromisso com uma concepção emancipadora de homens e mulheres que interagem no mundo e com o mundo.
Em uma perspectiva dialética, as DCNEI ( BRASIL, 2010a ) nos ajudam a entender esta ideia ao definir a concepção de \textbf{criança} como sujeito histórico, que ocupa lugar em um tempo real, e revela -se sujeito cultural como ser de criação e produtora de cultura.
A esse respeito, > considerar a \textbf{criança} como sujeito é levar em conta, nas relações que com ela estabelecemos que ela tem \textbf{desejos}, ideias, opiniões, capacidade de decidir, de criar, de \textbf{inventar}, que se manifestam desde cedo, nos seus movimentos, nas suas expressões, no seu olhar, nas suas vocalizações, na sua fala ( FARIA ; SALES, 2012, p. 56 ).
A relação estabelecida pelas autoras aproxima e intercruza os conceitos de sujeito e \textbf{criança}.
O resultado desta aproximação nos revela que a \textbf{criança} é um sujeito e é como sujeito que chega à escola, e não meramente como “ aluno ”.
Nesta condição de ser humano e de pessoa, a \textbf{criança} deve ser considerada em suas especificidades e linguagens, características múltiplas que devem ser conhecidas por todos os profissionais que com ela se relacionam.
É necessário que a concepção de criança-sujeito-histórico-cultural faça sentido na prática escolar, na vida dos professores, gestores, coordenadores pedagógicos e ganhe campo de atuação significativa na identidade para cada pessoa que é partícipe da comunidade escolar.
Considerando a concepção de criança-sujeito-histórico-cultural o \textbf{brincar} se mostra como expressão legítima onde se intercruzam múltiplas linguagens reveladas em pensamento e movimento que exercitam autonomia, argumentação, criação, direitos dentre tantas outras premissas constituintes do ser \textbf{criança}, articulados em expressão e linguagem e como ação promotora de aprendizagem.
Nesse sentido, o \textbf{brincar}, como experiência da \textbf{criança}, deve passar pelo crivo do \textbf{sensorial} das relações travadas e construídas entre os \textbf{adultos}, entre \textbf{adultos} e \textbf{crianças}, entre \textbf{crianças} e \textbf{crianças}.
As DCNEI ( BRASIL, 2010a ) estabelecem que as interações e a \textbf{brincadeira} são os eixos norteadores de práticas promotoras do aprender por meio de situações que efetivamente apresentem significado para a \textbf{criança} e ou grupo do qual faz parte.
Para tanto as experiências devem fazer sentido para elas nos contextos que falem e dialoguem sobre o mundo delas, mundo este do qual o \textbf{adulto} deve se aproximar para conhecer, interagir para que dele também possa aprender. \textbf{Brincar} implica em estabelecer vinculações entre o plano imaginário e o real, uma vez que a \textbf{criança} reproduz a realidade ao mesmo tempo que articula com o plano da imaginação.
Na Educação \textbf{Infantil}, no âmbito da escola, cabe ao professor proporcionar experiências ricas e diversificadas visando a > [... ] observar e constituir uma visão dos processos de desenvolvimento das \textbf{crianças} em conjunto e de cada uma em particular, registrando suas capacidades de uso das linguagens, assim como de suas capacidades sociais e dos recursos afetivos e emocionais que dispõem ( BRASIL, 1998c, 28 ).
As DCNEI ( BRASIL, 2010a ) enfatizam o direito da \textbf{criança} de viver a infância e de se desenvolver e destacam a Educação \textbf{Infantil} como o lugar do encontro em que as experiências acontecerão de modo que, por meio dessas vivências as \textbf{crianças} poderão amadurecer suas compreensões acerca da vida e do mundo, de si mesmas e do “ outro ”.
Paralelo a isso, colocam em prática “ formas de agir, sentir e pensar ” ( BRASIL, 2010a, p. 93 ).
É interessante destacar que o \textbf{brincar} não só existe e tem sentido no momento em que a \textbf{criança} chega à Escola ou à Educação \textbf{Infantil}, o \textbf{brincar} e a \textbf{brincadeira} traduzem e revelam quem são as \textbf{crianças}, como pensam ou organizam seus pensamentos, o que vivenciam em seu cotidiano e culturas, bem como revelam suas interações com \textbf{adultos} e seus pares, portanto, este espaço não é suficiente para abarcar sua abrangência.
É preciso compreender como nos diferentes tempos e culturas se pensou o \textbf{brincar} e o modo como às concepções foram se reconfigurando e até mesmo se equivocando a partir do momento que adentraram no espaço escolar.

% matched lemmas: adulto, brincadeira, brincar, criança, desejo, infantil, inventar, sensorial
\end{textsample}
